\subsection{同态的分解}

命题6. 设G是一个带算子的群，G'是一个带有Ω作用的原群，作用以指数形式书写。设f: G → G'是原群G到原群G'的一个同态，使得对所有α \ensuremath{\in} Ω及所有x \ensuremath{\in} G， \(f(x^{\alpha}) = f(x)^{\alpha}\) 。则 \(f(\mathbf{G})\) 是G'在G'上的运算及Ω作用下的稳定子集；该集合 \(f(\mathbf{G})\) 在诱导运算下是一个带算子的群，且该映射 \(x \mapsto f(x)\) 是一个群同态，其核为 G 的中心，其像为 \(f(\mathbf{G})\) 是带算子群的同态。

根据§1第4号命题1， \(f(\mathbf{G})\) 是……的稳定子集 \(\mathbf{G}'\) 根据内部法律规定， \(\mathbf{G}'\) 。对于每个元素 \(x \in \mathbf{G}\) 以及对于每个算子 \(\alpha\) , \(f(x)^{\alpha} = f(x^{\alpha}) \in f(\mathbf{G})\) ，从而 \(f(\mathbf{G})\) 在 \(\Omega\) 在 \(\mathbf{G}'\) 的作用下是稳定的。将 \(\mathbf{G}'\) 的内部运算写成乘法形式，

\[
\begin{array}{r l} (f (x) f (y)) f (z) & = f (x y) f (z) = f ((x y) z) = f (x (y z)) = f (x) f (y z) \\ & \qquad \qquad \qquad \qquad \qquad \qquad \qquad \qquad \qquad \qquad = f (x) (f (y) f (z)) \end{array}
\]

对于所有元素 \(x, y, z\) 在 \(\mathbf{G}\) ；因此诱导的定律在 \(f(\mathbf{G})\) 是结合的。

设 \(e\) 是 \(G\) 的单位元。其像 \(f(e)\) 是 \(f(G)\) 的单位元（§ 2，第 1 号）。的每个元素 \(f(G)\) 在 \(f(G)\) (§ 2，第 3 号)。因此，在……上诱导的法则 \(f(G)\) 根据内部法律规定， \(G'\) 是一个群法则。对于所有元素 \(x\) 和 \(y\) 在 \(G\) 以及每个算子 \(\alpha\) ,

\[
(f (x) f (y)) ^ {\alpha} = (f (x y)) ^ {\alpha} = f ((x y) ^ {\alpha}) = f \left(x ^ {\alpha} y ^ {\alpha}\right) = f \left(x ^ {\alpha}\right) f \left(y ^ {\alpha}\right) = (f (x)) ^ {\alpha} (f (y)) ^ {\alpha}
\]

这表明……的作用 \(\Omega\) 对群运算具有分配性 \(f(\mathrm{G})\) 的单位同态。因此 \(f(\mathrm{G})\) 在诱导运算下是一个带算子的群，并且显然该映射 \(x \mapsto f(x)\) 是带算子群的同态。

定义8. 设 \(f: \mathbf{G} \to \mathbf{G}'\) 设为单位元的逆像。 \(\mathbf{G}'\) 被称为 \(f\) .

的核 \(f\) 通常记作 \(\operatorname{Ker}(f)\) 以及图像 \(f(\mathbf{G})\) 于 \(f\) 有时用……表示 \(\operatorname{Im}(f)\) .

定理 3. 设 \(f: \mathbf{G} \to \mathbf{G}'\) 设为一个带算子的群同态。

\begin{enumerate}
\def\labelenumi{(\alph{enumi})}
\item
  \(\operatorname{Ker}(f)\) 是 \(\mathbf{G}\) ;
\item
  \(\operatorname{Im}(f)\) 是 \(\mathbf{G}'\) ;
\item
  映射 \(f\) 与定义在……上的等价关系相容 \(\mathbf{G}\) 中条件的扩张 \(\operatorname{Ker}(f)\)
\item
  映射 \(\tilde{f}\colon \mathbf{G} / \mathbf{Ker}(f)\to \operatorname {Im}(f)\) 通过取商导出的映射 \(f\) 通过取商，这是带算子的群的同构；
\item
  \(f = \iota \circ \tilde{f} \circ \pi\) 的x和y，其中 \(\iota\) 是以下映射的典范嵌入 \(\operatorname{Im}(f)\) 到 \(G'\) 和 \(\pi\) 是 \(G\) 到 \(G / \operatorname{Ker}(f)\) .
\end{enumerate}

的典范同态。断言(b)由命题6得出。等价关系 \(f(x) = f(y)\) 在 \(G\) 与 \(G\) 上的带算子群结构相容。由定理2（第4节），它因此具有形式 \(y \in xH\) 的x和y，其中 \(H\) 是 \(G\) 和 \(H\) 是单位元的类，因此 \(H = \operatorname{Ker}(f)\) 。于是断言(a)、(c)和(d)成立。断言(e)是显然的（集合论，II，§6，第5节）。

